\chapter{Теоретическая часть}

\section*{Что собой представляет программа <<Телефонный справочник>> на Prolog?}

Программа на языке Prolog --- база знаний и вопрос, на который система должна дать ответ и то, как она к этому пришла.
База знаний представляет собой набор фактов с правилами.
В данном  случае база знаний хранит информацию <<Телефонный справочник>>: Фамилия, номер телефона, адрес --- структура (город, улица, номер дома, номер квартиры) и <<Автомобили>>: Фамилия, марка, цвет, стоимость, номер.

\section*{Какова ее структура?}

Программа на Prolog состоит из разделов. Каждый раздел начинается со своего заголовка. 
В программе не обязательно должны быть все разделы.
Структура программы:
\begin{itemize}
	\item директивы компилятора --- раздел описания констант;
	\item DOMAINS --- раздел описания доменов;
	\item DATABASE --- раздел описания предикатов внутренней базы данных;
	\item PREDICATES --- раздел описания предикатов;
	\item CLAUSES --- раздел описания предложений базы знаний;
	\item GOAL ---  раздел описания внутренней цели (вопроса).
\end{itemize}

В данной лаб. раб. используются CLAUSES, GOAL и DOMAINS.

\section*{Как она реализуется, как формируются результаты работы программы?}

Prolog включает механизм вывода, который основан на сопоставлении образцов.
Программа на Prolog представляет собой: базу знаний и вопрос. С помощью подбора ответов
на запросы он (программа) извлекает хранящуюся (известную в программе) информацию.
База знаний содержит истинностные знания, используя которые программа выдает ответ на
запрос. Особенностью Prolog является то, что при поиске ответов на вопрос, он
рассматривает альтернативные варианты и находит все возможные решения (методом проб и
ошибок) --- множества значений переменных, при которых на поставленный вопрос можно
ответить --- <<да>>.
